\section{RSP: de la Historical ASL a las Capability-Triggered Safeguards}

En la clase se usan los AI Safety Levels para describir la relación entre la capacidad de un modelo y su nivel de protección, y se enfatiza la defense in depth. La versión actual, Anthropic RSP v3.0, evolucionó hacia capability thresholds, required safeguards, risk reports y procesos de gobernanza. La durable idea que comparten ambas versiones es esta: a medida que la capability evidence se acerca a un threshold de alto impacto, el entrenamiento y el despliegue deben elevar primero los security/deployment controls, en lugar de corregir después.

\subsection{Threshold, safeguard y gate}

Primero, el equipo reúne capability evidence mediante evaluation, elicitation y forecast, y después determina qué tan lejos está del threshold. Si lo alcanza o se acerca a él, debe cumplir las required safeguards correspondientes, por ejemplo una model-weight security más sólida, access control, runtime guard, monitoring, red team, incident readiness y governance review. Solo al final decide si continúa el entrenamiento, amplía el despliegue o restringe el access.

\lecturefigure{12-capability-safeguards.png}{La evidencia de capacidad activa la threshold review, las required safeguards y el deployment gate.}{ASL framing de la clase; Anthropic RSP v3.0; redibujo conceptual.}

\begin{knowledgebox}{Lectura de la figura: el cambio de versión muestra que la governance debe poder revisarse}
La ASL de la clase vinculaba de forma bastante estrecha los niveles del modelo y los de protección; las versiones posteriores de la RSP separan el capability threshold del safeguard package, para poder actualizarlos en función de riesgos específicos. Este cambio no es una simple sustitución del tipo «la versión anterior estaba mal, la nueva está bien», sino el reconocimiento de que la evaluation science, el threat model y la safeguard evidence cambian; por lo tanto, la policy misma también debe estar versioned, publicar sus cambios y someterse a review.
\end{knowledgebox}

\teachervoice{El ponente enfatiza que una de las mayores dificultades de la evaluación es el elicitation overhang: que hoy no se haya detectado una capacidad no significa que siga sin lograrse con un scaffold, una tool o un research effort más potentes. En la clase, la ASL no es una etiqueta permanente que se le pone al modelo, sino una exigencia de que la organización prepare las protecciones por adelantado en condiciones de incertidumbre.}

\subsection{Defense in Depth}

La \term{defense in depth} (defensa en profundidad) consiste en establecer varias capas de control independientes entre sí y con modos de falla distintos. Training reduce las conductas peligrosas mediante data/objective y post-training; la capa de access limita quién puede usar qué capacidad; el runtime classifier/tool policy supervisa las solicitudes y las action; la capa de response se encarga del red team, incident, patch, revoke y la customer communication.

\lecturefigure{13-defense-in-depth.png}{La seguridad de un frontier model requiere defensa en profundidad en training, access, runtime y response.}{Subtítulos locales 37:00--39:10; redibujo conceptual.}

\begin{knowledgebox}{Lectura de la figura: más capas no implica independencia}
Si training, el classifier y el judge dependen todos del mismo modelo y de los mismos datos, pueden compartir un blind spot. Una defensa en profundidad eficaz requiere analizar la common-mode failure, la separación de permisos, la independencia de la supervisión y el human override. Cada capa debe producir evidence para que la siguiente capa juzgue, y contar con un fallback explícito, en lugar de reenviar el riesgo en círculo de una capa a otra.
\end{knowledgebox}

\subsection{Resumen de la sección}

El núcleo de la RSP es la capability-triggered governance. Los términos pueden cambiar; el durable contract es: primero evaluar, luego cumplir las required safeguards y por último entrenar/desplegar; toda excepción debe tener un owner, evidencia y una fecha de reevaluación.
